\documentclass[11pt]{article}
\usepackage[margin=1in]{geometry}
\usepackage{enumitem}
% =============================================================================
% Preambles/header.tex — shared LaTeX/Beamer preamble for this project
%
% Use from a Beamer lecture by prepending:
%     \input{header}
% and compiling with TEXINPUTS=../Preambles:$TEXINPUTS (the /compile-latex
% skill does this automatically).
%
% PALETTE CONTRACT. Color names defined here MUST match the names in
% Quarto/theme-template.scss so Beamer and Quarto renderings use the same
% palette. The scripts/check-palette-sync.sh script enforces this.
%
% To customize: edit the HEX values below AND the matching $variable in
% Quarto/theme-template.scss, then run scripts/check-palette-sync.sh.
% =============================================================================

% -----------------------------------------------------------------------------
% Engine detection + encoding
%
% Our /compile-latex skill uses XeLaTeX, where inputenc is unnecessary and
% can produce encoding warnings. Only load inputenc under pdfTeX.
% -----------------------------------------------------------------------------
\usepackage{iftex}
\ifPDFTeX
  \usepackage[utf8]{inputenc}
\fi

\usepackage{amsmath, amssymb, amsthm}
\usepackage{booktabs}
\usepackage{graphicx}

% -----------------------------------------------------------------------------
% PALETTE — mirrors Quarto/theme-template.scss variable names.
% Load xcolor and define colors BEFORE anything that references them
% (notably hyperref, hypersetup, and the Beamer color assignments).
% -----------------------------------------------------------------------------
\usepackage{xcolor}

\definecolor{primary-blue}{HTML}{012169}      % headings, accents
\definecolor{primary-gold}{HTML}{B9975B}      % emphasis, borders
\definecolor{highlight-yellow}{HTML}{F2A900}  % markers, alerts
\definecolor{light-bg}{HTML}{E8EDF5}          % title / section backgrounds
\definecolor{jet}{HTML}{1A1A1A}               % body text

% Semantic colors (used in diagrams and annotations)
\definecolor{positive}{HTML}{15803D}          % good / observed / identified
\definecolor{negative}{HTML}{B91C1C}          % bad / problematic / bias
\definecolor{neutral}{HTML}{525252}           % reference / context

% Highlight accents (match .hi-* classes in the SCSS)
\definecolor{hi-slate}{HTML}{314F4F}
\definecolor{hi-green}{HTML}{2E8B57}
\definecolor{hi-red}{HTML}{C41E3A}

% Semantic aliases so skill templates and snippets can write
% \textcolor{accent}{...} without hard-coding "primary-blue".
\colorlet{accent}{primary-blue}
\colorlet{accent2}{primary-gold}

% -----------------------------------------------------------------------------
% Hyperref — loaded AFTER xcolor + palette so link colors resolve.
% -----------------------------------------------------------------------------
\usepackage{hyperref}
\hypersetup{colorlinks=true, linkcolor=primary-blue, urlcolor=primary-blue,
            citecolor=primary-blue, pdfborder={0 0 0}}

% -----------------------------------------------------------------------------
% Course code — set once per deck (after \input{header}, before \begin{document})
% via \coursecode{CS401}. Shown in the footer on every slide so a deck's course
% is identifiable without opening the title page. Empty by default.
% -----------------------------------------------------------------------------
\makeatletter
\newcommand{\coursecode}[1]{\gdef\@coursecodetext{#1}}
\gdef\@coursecodetext{}
\makeatother

% -----------------------------------------------------------------------------
% Beamer theme assignments (only apply under Beamer).
%
% We detect Beamer via \@ifclassloaded{beamer}, which is the documented,
% non-internal way. This must be wrapped in \makeatletter/\makeatother
% because \@ifclassloaded contains an @-catcode character.
% -----------------------------------------------------------------------------
\makeatletter
\@ifclassloaded{beamer}{%
  \usefonttheme{professionalfonts}
  \setbeamercolor{structure}{fg=primary-blue}
  \setbeamercolor{frametitle}{fg=primary-blue}
  \setbeamercolor{title}{fg=primary-blue}
  \setbeamercolor{subtitle}{fg=primary-gold}
  \setbeamercolor{author}{fg=primary-blue}
  \setbeamercolor{institute}{fg=neutral}
  \setbeamercolor{date}{fg=primary-gold}
  \setbeamercolor{itemize item}{fg=primary-blue}
  \setbeamercolor{itemize subitem}{fg=primary-gold}
  \setbeamercolor{alerted text}{fg=primary-gold}
  \setbeamercolor{block title}{fg=white, bg=primary-blue}
  \setbeamercolor{block body}{bg=light-bg}
  % Semantic block variants: block=definition/concept (blue), exampleblock=
  % worked example (green/positive), alertblock=key takeaway or caution (gold).
  % Keep to <=2 colored boxes per slide across ALL three variants (INV-7).
  \setbeamercolor{block title example}{fg=white, bg=positive}
  \setbeamercolor{block body example}{bg=light-bg}
  \setbeamercolor{block title alerted}{fg=white, bg=primary-gold}
  \setbeamercolor{block body alerted}{bg=light-bg}

  % Minimal footer: course code (left) + page X / N (right), no navigation clutter
  \setbeamertemplate{navigation symbols}{}
  \setbeamertemplate{footline}{%
    \color{neutral}\footnotesize\hspace{0.7em}\@coursecodetext\hfill
    \insertframenumber/\inserttotalframenumber\hspace{0.7em}\vspace{0.3em}}
}{}
\makeatother

% -----------------------------------------------------------------------------
% TikZ setup — libraries the snippet gallery uses
% -----------------------------------------------------------------------------
\usepackage{tikz}
\usetikzlibrary{arrows.meta, positioning, calc, decorations.pathreplacing,
                fit, shapes.geometric, backgrounds}

% Shared TikZ styles — snippets and hand-written diagrams can pick these up
% via `\begin{tikzpicture}[every node/.style={...}]` or by naming specific
% styles (e.g., `\node[dag-node] (X) at (0,0) {X};`).
\tikzset{
  % Node styles for causal DAGs and similar
  dag-node/.style = {
    draw=primary-blue, fill=light-bg, rounded corners=2pt,
    minimum width=1.2cm, minimum height=0.9cm, align=center, font=\small
  },
  decision-node/.style = {
    draw=primary-gold, fill=white, diamond, aspect=1.8,
    minimum width=1.8cm, minimum height=1.2cm, align=center, font=\small,
    inner sep=1pt
  },
  % Edge styles — observed vs counterfactual
  observed-edge/.style = {-{Latex[length=2.5mm]}, thick, draw=jet},
  counterfactual-edge/.style = {-{Latex[length=2.5mm]}, thick, draw=neutral, dashed},
  confound-edge/.style = {-{Latex[length=2.5mm]}, thick, draw=negative, dashed},
  % Data-point styles for plots
  observed-dot/.style = {fill=primary-blue, draw=primary-blue, circle, inner sep=1.2pt},
  counterfactual-dot/.style = {fill=white, draw=neutral, circle, inner sep=1.2pt}
}

% -----------------------------------------------------------------------------
% Convenience macros
% -----------------------------------------------------------------------------
\newcommand{\muted}[1]{\textcolor{neutral}{#1}}
\newcommand{\key}[1]{\textcolor{primary-gold}{\textbf{#1}}}
\newcommand{\good}[1]{\textcolor{positive}{#1}}
\newcommand{\bad}[1]{\textcolor{negative}{#1}}

% Standout frame macro: full-bleed dark background for section transitions.
% Beamer-only (uses \begin{frame}/\setbeamercolor/\usebeamerfont) -- do not
% call from an article-class file (e.g. Notes/<CODE>/*-notes.tex). \muted,
% \key, \good, \bad, and \coursecode above are plain xcolor/gdef and are
% safe to use from either class; verified when Notes/article-class support
% was added.
% Expand: \transitionslide{Your transition title}
\newcommand{\transitionslide}[1]{%
  \begingroup
  \setbeamercolor{background canvas}{bg=primary-blue}
  \begin{frame}[plain]
    \centering
    \vfill
    {\usebeamerfont{title}\color{white}\Large #1\par}
    \vfill
  \end{frame}
  \endgroup
}

% Section divider frame (Beamer-only), an alternative to \transitionslide with a
% darker two-line style: near-black (jet) background, a small white label line
% above a larger gold title, and a thin gold rule along the bottom edge.
% Expand: \sectiondivider{Section 2}{Pointers}
\newcommand{\sectiondivider}[2]{%
  \begingroup
  \setbeamercolor{background canvas}{bg=jet}
  \begin{frame}[plain]
    \vspace*{3.0cm}
    {\color{white}\ttfamily\large #1\par}
    \vspace{0.5cm}
    {\color{primary-gold}\LARGE #2\par}
    \begin{tikzpicture}[remember picture, overlay]
      \draw[primary-gold, line width=2.5pt]
        ($(current page.south west)+(0,0.1)$) -- ($(current page.south east)+(0,0.1)$);
    \end{tikzpicture}
  \end{frame}
  \endgroup
}

% -----------------------------------------------------------------------------
% Channel watermark — "Concepts That Click" (YouTube).
%
% Article-class files (CompetitiveExam Q/A sets, Notes, Assignments) get a
% light diagonal draft watermark on every page plus a centered footer naming
% the channel. Beamer decks get a subtle TikZ background watermark instead
% (the footline already carries the course code). Centralized here so every
% MCQ PDF is branded without per-file edits.
% -----------------------------------------------------------------------------
\makeatletter
\@ifclassloaded{article}{%
  \usepackage{draftwatermark}
  \SetWatermarkText{Concepts That Click}
  \SetWatermarkScale{1}
  \SetWatermarkFontSize{2cm}
  \SetWatermarkLightness{0.86}
  \SetWatermarkAngle{45}
  \usepackage{fancyhdr}
  \pagestyle{fancy}
  \fancyhf{}
  \fancyfoot[C]{\footnotesize\textcolor{neutral}{Concepts That Click~|~YouTube: \href{https://www.youtube.com/@concepts.thatclick}{@concepts.thatclick}}}
  \renewcommand{\headrulewidth}{0pt}
  \renewcommand{\footrulewidth}{0pt}
  \fancypagestyle{plain}{%
    \fancyhf{}%
    \fancyfoot[C]{\footnotesize\textcolor{neutral}{Concepts That Click~|~YouTube: \href{https://www.youtube.com/@concepts.thatclick}{@concepts.thatclick}}}%
    \renewcommand{\headrulewidth}{0pt}%
    \renewcommand{\footrulewidth}{0pt}%
  }
}{}
\@ifclassloaded{beamer}{%
  \setbeamertemplate{background}{%
    \begin{tikzpicture}[remember picture, overlay]
      \node[rotate=45, opacity=0.055, align=center]
        at (current page.center)
        {\Huge\textcolor{neutral}{Concepts That Click}};
    \end{tikzpicture}%
  }
}{}
\makeatother 
\coursecode{JKSSB}
\renewcommand{\thesection}{59.\arabic{section}}
\newtheorem{example}{Example}[section]
\newenvironment{definitionbox}[1]{%
  \par\noindent\textbf{Definition (#1).}\ \itshape
}{\par\vspace{0.3em}}
\title{Current CPU Architecture, Data Representation \& RISC/CISC --- Detailed Lecture Notes}
\author{JKSSB Computer Awareness}
\date{\today}

\begin{document}
\maketitle

\section{von Neumann architecture and the bottleneck}
A \key{von Neumann} (stored-program) computer keeps both instructions and
data in the same memory. The processor fetches an instruction, decodes it,
and executes it; when the instruction needs a data operand, that operand is
read from the same memory. Because a single shared path carries instructions
and data, the processor frequently has to wait for memory. This limitation is
called the \key{von Neumann bottleneck}: the rate at which work can proceed is
capped by the rate at which instructions and data can be transferred between
the CPU and memory over that shared path. Faster memory, a cache, and wider
paths all reduce --- but do not remove --- the bottleneck.

\begin{definitionbox}{von Neumann architecture}
A stored-program design in which instructions and data are held together in
the same memory and fetched over a shared path between the CPU and memory.
\end{definitionbox}

The important contrast is \emph{same memory for instructions and data}. A
system that keeps them physically separate is a different (Harvard-style)
design and is not the von Neumann model.

\section{Addressing modes: immediate addressing}
An \key{addressing mode} describes how an instruction states where its
operand is. \key{Immediate addressing} embeds the operand value directly
inside the instruction. For example, an instruction ``add 5 to the register''
carries the value 5 in the instruction encoding itself. Because the value is
already present, no extra memory access is required to obtain the operand.

The near-neighbour contrasts matter for exams:
\begin{itemize}
  \item \emph{Immediate}: the operand is in the instruction.
  \item \emph{Register}: the operand is in a named register.
  \item \emph{Direct}: the instruction names a memory address holding the
    operand.
\end{itemize}
The single line to remember is: immediate addressing means the operand is a
constant carried with the instruction.

\section{Cache and the cache-hit ratio}
A \key{cache} is a small, fast memory placed between the CPU and main memory.
It holds copies of recently used instructions and data. When the CPU asks for
an item and the cache has it, that access is a \key{hit}; if the cache does not
have it, the access is a \key{miss} and the item must come from main memory
(and is usually placed in the cache for next time).

\begin{definitionbox}{Cache-hit ratio}
\[
\text{hit ratio} = \frac{\text{number of hits}}{\text{total accesses}}
= \frac{\text{hits}}{\text{hits} + \text{misses}}.
\]
\end{definitionbox}

A \key{hit} is faster because cache access time is much lower than the access
time of main memory. On a hit the CPU avoids the slow main-memory access; on
a miss it pays that cost. This is why a higher hit ratio improves average
performance.

\begin{example}
If a program makes 900 cache hits and 100 cache misses, the total accesses
are $900+100=1000$ and the hit ratio is $900/1000 = 0.90$ (90\%).
\end{example}

\section{Data hazards even with register addressing}
A \key{data hazard} is a dependency that prevents an instruction from
proceeding correctly in the cycle originally planned. A common case is a
read-after-write dependency: an instruction needs the value of a register
that an earlier instruction has not yet written. Pipelining is what exposes
these dependencies, because several instructions are in progress at once.

Register addressing avoids an extra memory access and removes some hazards
caused by memory operands, but it does \bad{not} eliminate all hazards. Two
instructions can still depend on the same register. This is the trap: a
candidate may reason that ``no memory operand means no hazard,'' which is
false. Hazards arise from dependencies between instructions, whatever kind of
operand is used.

\section{Virtual memory versus cache}
Cache and virtual memory are frequently confused, but they solve different
problems.

\begin{center}
\begin{tabular}{p{0.20\textwidth}p{0.36\textwidth}p{0.36\textwidth}}
\toprule
\textbf{Point} & \textbf{Cache} & \textbf{Virtual memory} \\
\midrule
Location & Small fast memory next to the CPU. & Uses main memory plus secondary storage as backing store. \\
Purpose & Reduce average memory access time by keeping recently used blocks close. & Make the memory available to programs appear larger than physical RAM. \\
Management & Largely hardware. & Operating system with hardware support. \\
\bottomrule
\end{tabular}
\end{center}

Both involve small ``fast'' storage for larger ``slow'' storage, but the cache
shortens access time while virtual memory expands the apparent size of memory.
They are not the same mechanism and are not interchangeable.

\section{Instruction cycle: Program Counter and Instruction Register}
During the fetch stage the CPU reads an instruction from memory. The
\key{Program Counter (PC)} holds the address of the \key{next} instruction to
be fetched; once the current instruction has been fetched, the PC is advanced
to point at the following one. The \key{Instruction Register (IR)} holds the
instruction that is currently being fetched, decoded, and executed.

The exam trap is the PC. It is easy to say the PC holds ``the address of the
current instruction,'' but at the moment the instruction is executing, the PC
already points at the \emph{next} instruction. The correct statement is: the
PC holds the address of the next instruction to be fetched.

\section{Interrupts are processed at instruction boundaries}
An \key{interrupt} is a signal that causes the CPU to suspend the current
program temporarily and transfer control to a handler routine, after which
execution usually resumes. Interrupts are recognised at \key{instruction
boundaries} --- between one instruction and the next --- during the program's
execution, not only when the whole program finishes.

The trap is the claim that ``an interrupt is processed only after the entire
program completes.'' That is false. Mid-program interrupts are the normal
case: the processor finishes the instruction in progress, saves its state, and
services the interrupt before continuing.

\section{Instruction timing: clock cycles vary}
Different machine instructions perform different amounts of work. A simple
register-to-register operation may complete in a single clock cycle, while a
division, a complex operation, or an instruction that must access slow memory
can take many cycles. Instruction timing therefore depends on the
instruction, its operands, and the memory system.

The trap is the claim that ``all machine instructions take the same number of
clock cycles.'' The count varies by instruction; averaging over a program
gives an average cycles-per-instruction, not a constant per instruction.

\section{Two's complement and overflow detection}
In an $n$-bit \key{two's-complement} representation the most significant bit
is the \key{sign} bit: 0 for non-negative and 1 for negative values. A
negative number is formed by inverting the bits of its magnitude and adding
one. With $n$ bits the representable range is
$-2^{\,n-1}$ through $+2^{\,n-1}-1$.

\key{Overflow} occurs when the arithmetic result cannot be represented in the
available bits. The reliable test is the \key{sign} rule:

\begin{definitionbox}{Overflow test}
Overflow occurs when two operands of the \key{same sign} are added and the
result has the \key{opposite sign}. If the two operands have different signs,
no overflow occurs.
\end{definitionbox}

Overflow is \bad{not} the same as a carry out of the most significant bit. A
carry can be produced without overflow, and overflow can occur without the
final carry telling the full story. Use the sign test, not the carry flag.

\begin{example}[Overflow]
Work in 8-bit two's complement. Add $+7$ and $+1$:
\[
0111\ (+\!7) \;+\; 0001\ (+\!1) \;=\; 1000.
\]
The operands are both positive, but the sum $1000$ has sign bit 1 and
represents $-8$. Two same-sign operands produced an opposite-sign result, so
this is \good{overflow}.
\end{example}

\begin{example}[Carry without overflow]
Again in 8 bits, add $-1$ and $-1$:
\[
1111\,1111\ (-\!1) \;+\; 1111\,1111\ (-\!1) \;=\; 1\,1111\,1110.
\]
There is a carry out of the most significant bit, but the discarded 9th bit
leaves $1111\,1110$, which is exactly $-2$. The result is correct, so there is
\good{no overflow}. This shows that a carry out of the MSB is not the same as
overflow.
\end{example}

\section{Zero-extension versus sign-extension}
\key{Extension} widens a value from fewer bits to more bits while preserving
its meaning. The rule depends on whether the value is unsigned or signed:
zero-extension fills the new leftmost bits with 0 and is used for
\key{unsigned} values; sign-extension copies the sign bit into the new
leftmost bits and is used for \key{signed} values. The common trap is to swap
the two.

\begin{center}
\begin{tabular}{llll}
\toprule
\textbf{8-bit pattern} & \textbf{Unsigned} & \textbf{Signed} & \textbf{16-bit sign-extended} \\
\midrule
$0000\,0101$ & $5$ & $+5$ & $0000\,0000\,0000\,0101$ \\
$1111\,1011$ & $251$ & $-5$ & $1111\,1111\,1111\,1011$ \\
\bottomrule
\end{tabular}
\end{center}

Zero-extending $1111\,1011$ would give $0000\,0000\,1111\,1011$ (value $251$),
which preserves the \emph{unsigned} value. Sign-extending it gives
$1111\,1111\,1111\,1011$ (value $-5$), which preserves the \emph{signed}
value. Using the wrong rule changes the value.

\section{BCD: binary-coded decimal}
\key{BCD} (Binary-Coded Decimal) represents each \key{decimal digit} with its
own group of \key{4 bits}. The ten digit codes are:

\begin{center}
\begin{tabular}{cccccccccc}
\toprule
$0$ & $1$ & $2$ & $3$ & $4$ & $5$ & $6$ & $7$ & $8$ & $9$ \\
\midrule
$0000$ & $0001$ & $0010$ & $0011$ & $0100$ & $0101$ & $0110$ & $0111$ & $1000$ & $1001$ \\
\bottomrule
\end{tabular}
\end{center}

\begin{example}[BCD encoding]
Encode the decimal number $59$ in BCD. Take each digit separately: $5$ is
$0101$ and $9$ is $1001$. So $59$ becomes
\[
0101\ 1001 \quad\text{(BCD)}.
\]
For contrast, the ordinary binary value of $59$ is $111011_2$. BCD keeps one
4-bit code per decimal digit; it does not pack the whole number as plain
binary.
\end{example}

The trap is the claim that BCD is ``just the binary of the whole number.''
It is not: each decimal digit is coded in 4 bits, and the codes $1010$ through
$1111$ are not valid BCD digits.

\section{Byte-addressability and effective address vs instruction length}
Memory is \key{byte-addressable} when every individual byte has its own unique
address. A word is a group of bytes at consecutive addresses. Because bytes
are individually addressable, a word access begins at the address of its
first byte.

The \key{effective address} is the actual location that an instruction uses
for its operand. Depending on the addressing mode, it may be given directly
or computed at run time from registers and constants. Instruction length is
the number of bytes the instruction itself occupies. These are different
quantities: a long instruction does not imply a distant operand. The address
of the operand is set by the addressing mode, not by the instruction's size.
Confusing the two is the trap; keep ``how big is the instruction'' separate
from ``where is the operand.''

\section{RISC versus CISC}
The two broad instruction-set philosophies are \key{RISC} (Reduced Instruction
Set Computer) and \key{CISC} (Complex Instruction Set Computer).

\begin{center}
\begin{tabular}{p{0.22\textwidth}p{0.34\textwidth}p{0.34\textwidth}}
\toprule
\textbf{Aspect} & \textbf{RISC} & \textbf{CISC} \\
\midrule
Instruction set & Small, simple, regular. & Large and complex. \\
Instruction length & Fixed. & Variable. \\
Registers & Many general-purpose registers. & Fewer registers. \\
Memory access & Load/store architecture: only load and store touch memory. & Instructions may operate directly on memory operands. \\
Cycles per instruction & Typically one, suited to pipelining. & Often multiple cycles. \\
Control emphasis & Hardwired control; work shifted to the compiler. & Microcoded control; work done in hardware. \\
Typical examples & ARM, MIPS, RISC-V. & x86 and older mainframe designs. \\
\bottomrule
\end{tabular}
\end{center}

RISC keeps instructions simple and fixed-length so they pipeline well, relying
on many registers and a load/store discipline. CISC offers powerful,
variable-length instructions that can do more in one step, at the cost of
more complex decoding. The exam cue is the pair: \emph{RISC = reduced, fixed,
many registers, load/store}; \emph{CISC = complex, variable-length, memory
operands allowed}.

\subsection*{Exercises}
\begin{enumerate}[label=\arabic*.]
  \item State what is shared in von Neumann architecture and define the von
    Neumann bottleneck.
  \item What is immediate addressing, and why does it need no extra memory
    access for the operand?
  \item A program makes 850 cache hits and 150 misses. Compute the
    cache-hit ratio.
  \item Explain why register addressing does not eliminate all data hazards.
  \item In 8-bit two's complement, add $+7$ and $+1$ and state whether
    overflow occurs. Justify using the sign test.
  \item Encode the decimal number $59$ in BCD, and contrast it with the plain
    binary value.
\end{enumerate}

\subsection*{Solutions}
\begingroup\small
\begin{enumerate}[label=\arabic*.]
  \item Instructions and data share the same memory. The bottleneck is the
    limited rate at which instructions and data can be transferred between
    the CPU and memory over the shared path.
  \item Immediate addressing carries the operand inside the instruction
    encoding, so the value is already available and no memory fetch is
    required for it.
  \item Total accesses $=850+150=1000$, so the hit ratio is
    $850/1000 = 0.85$ (85\%).
  \item Register operands avoid a memory access, but two instructions can
    still depend on the same register (for example a read-after-write
    dependency), so hazards remain.
  \item $0111\ (+\!7) + 0001\ (+\!1) = 1000$, which has sign bit 1 and
    represents $-8$. Two same-sign operands produced an opposite-sign
    result, so overflow occurs.
  \item BCD: $5 \rightarrow 0101$, $9 \rightarrow 1001$, giving
    $0101\ 1001$. Plain binary of $59$ is $111011_2$; BCD keeps one 4-bit
    code per decimal digit instead of packing the number as binary.
\end{enumerate}
\endgroup
\end{document}
